% Generated draft: requires mathematical and source review.
\documentclass[11pt]{article}
\usepackage[margin=25mm]{geometry}
\usepackage{fontspec}
\setmainfont{Latin Modern Roman}
\usepackage{amsmath,amssymb,mathtools,unicode-math}
\setmathfont{Latin Modern Math}
\usepackage{enumitem,needspace,xurl,hyperref}
\setlength{\parindent}{0pt}
\setlength{\parskip}{5pt}
\setlength{\emergencystretch}{3em}
\newcommand{\N}{\mathbb{N}}
\newcommand{\Z}{\mathbb{Z}}
\newcommand{\Q}{\mathbb{Q}}
\newcommand{\R}{\mathbb{R}}
\newcommand{\C}{\mathbb{C}}
\newcommand{\F}{\mathbb{F}}
\newcommand{\A}{\mathbb{A}}
\newcommand{\PP}{\mathbb{P}}
\newcommand{\E}{\mathbb{E}}
\newcommand{\eps}{\varepsilon}
\newcommand{\dd}{\,\mathrm{d}}
\begin{document}
\section*{UnsolvedMath ID: 30000538}
\textbf{Source problem code:} catalogue entry
\subsection{Definitions and mathematical statement}
\textbf{Extracted \(R_\infty\) question.} Does every relatively hyperbolic group, and hence every Kleinian group, have property \(R_\infty\)? For \(\varphi\in\operatorname{Aut}(G)\), elements \(x,y\in G\) are \(\varphi\)-twisted conjugate if \(y=gx\varphi(g)^{-1}\) for some \(g\in G\). The Reidemeister number \(R(\varphi)\) is the number of these classes. Property \(R_\infty\) means \(R(\varphi)=\infty\) for every automorphism, not merely for one automorphism or for ordinary conjugacy.\par \textbf{Essential scope gap.} The literal all-groups wording includes elementary examples that do not have this property. For \(G=\Z\) and \(\varphi(n)=-n\), twisted classes are the two residue classes modulo \(2\), so \(R(\varphi)=2\). The group \(\Z\) is hyperbolic and occurs as an infinite cyclic Kleinian group generated by a loxodromic isometry. Thus an omitted non-elementarity or other qualifying restriction is essential to the intended research question; the unrestricted wording cannot simply be advertised as an unresolved universal assertion. This is a formulation diagnostic, not a claim to settle the intended qualified problem.\par The source does not specify the definition and peripheral structure of relative hyperbolicity, the required finiteness hypotheses, or whether automorphisms must preserve that structure. A Kleinian group conventionally means a discrete subgroup of \(\mathrm{PSL}_2(\C)\); the word alone does not supply the non-elementarity or finite-generation/geometric hypotheses needed for the stated \emph{and hence} passage to a selected relatively hyperbolic class. These restrictions are not guessed. Once the intended class \(\mathcal C\) is recovered, the actual target is \(\forall G\in\mathcal C\,\forall\varphi\in\operatorname{Aut}(G),\ R(\varphi)=\infty\), with its precise automorphism scope.\par The supplied triage checked 2026-08-21 labels the item partially solved and reports \(R_\infty\) for non-elementary hyperbolic groups and broad relatively hyperbolic subclasses, while calling an unrestricted relatively hyperbolic assertion open. The elementary obstruction above shows that this last phrase needs reconciliation with the intended missing hypotheses. The cited later theorem's exact scope is not reproduced and no independent current-status conclusion is drawn.
\subsection{Short English statement}
Property \(R_\infty\) requires infinitely many twisted conjugacy classes for every automorphism. The literal universal claim needs a missing restriction: \(\Z\), also an elementary Kleinian group, has only two classes for inversion. The intended non-elementary or peripheral hypotheses must be recovered.
\subsection{Sources}
\begin{itemize}\item Local UnsolvedMath registry, permanent ID 30000538, code catalogue entry; \emph{\(R_\infty\) Property for Relatively Hyperbolic Groups}, Oberwolfach Reports Open Problems.\item \emph{Geometric Group Theory, Hyperbolic Dynamics and Symplectic Geometry (2007)}, Oberwolfach Reports, DOI 10.4171/owr/2006/33, pp. 1991--2058. \url{https://doi.org/10.4171/owr/2006/33}. Proposed year and year tag: 2006; no proposer named.\item Supplied literature label \emph{OWR Geometric Group Theory report}: \url{https://doi.org/10.4171/owr/2006/33}.\item Supplied literature label \emph{Fournier-Facio--Iveson--Martino--Sgobbi--Wong}: \url{https://doi.org/10.1007/s10711-026-01078-x}.\item Literature triage checked 2026-08-21, status \texttt{partially\_solved}; record updated 2026-09-16. The \texttt{corrected\_verified} field claims editorial and notation cleanup only. Verify dated claims independently.\end{itemize}
\end{document}
